
\section*{MoE in Clinical Natural Language Processing and Medical LLMs}
\addcontentsline{toc}{section}{MoE in Clinical NLP and Medical LLMs}

The application of MoE architectures to clinical natural language processing (NLP) and medical large language models (LLMs) represents one of the most rapidly evolving frontiers in healthcare AI. Clinical text is inherently heterogeneous, encompassing diverse document types (discharge summaries, radiology reports, pathology notes, surgical records), medical specialties, and linguistic styles. MoE architectures offer a natural solution by enabling specialization of different experts for distinct clinical domains and document types.

\subsection*{Medical LLMs with MoE architecture}

The emergence of foundation models for medicine has been profoundly influenced by MoE design principles. Medical LLMs face a fundamental tension: they must possess broad medical knowledge spanning multiple specialties while maintaining sufficient depth in specific clinical domains. Dense models address this by uniformly scaling all parameters, but MoE architectures achieve this more efficiently by routing different medical queries to specialized expert pathways.

\cite{bechar2025biocancer_moe} introduced BioCancer-MoLA, a Mixture of LoRA Experts framework designed for breast cancer clinical decision support. The architecture employs domain-specific low-rank adaptation (LoRA) experts, each fine-tuned on a different aspect of breast cancer knowledge (pathology, radiology, oncology, genetics), with a knowledge-guided routing mechanism that directs clinical queries to the most relevant experts. Critically, BioCancer-MoLA incorporates a hallucination mitigation strategy that uses domain knowledge consistency checks to filter unreliable expert outputs, achieving a 34\% reduction in clinically significant hallucinations compared to monolithic fine-tuned models. This work highlights a crucial insight: in high-stakes medical applications, MoE not only improves performance but can also enhance safety by enabling targeted quality control at the expert level.

\cite{chen2026cor_mole} proposed CoR-MoLE (Specialist-Collaborative Reasoning with Mixture-of-LoRA-Experts), which simulates multi-disciplinary team (MDT) decision-making through collaborative expert reasoning. In clinical oncology, MDT discussions are the gold standard for complex case management, where specialists from different disciplines (surgery, medical oncology, radiation oncology, pathology, radiology) contribute their expertise. CoR-MoLE mirrors this structure by training discipline-specific expert adapters and implementing a collaborative reasoning protocol where experts exchange intermediate assessments before reaching a consensus recommendation. This approach achieved superior diagnostic accuracy compared to single-expert and ensemble baselines on complex oncology cases, demonstrating that MoE architectures can capture not just individual expertise but also the dynamics of clinical collaboration.

\subsection*{Parameter-efficient MoE adaptation}

Training or fine-tuning large medical LLMs with full MoE architectures remains computationally expensive. The combination of MoE with parameter-efficient fine-tuning (PEFT) methods has emerged as a practical solution. \cite{wang2025doublegate} introduced DoubleGate, a dual-gating Mixture-of-LoRA-Experts framework for medical multi-task learning. DoubleGate maintains two levels of gating: a task-level gate that selects which expert pool to consult, and an expert-level gate that performs fine-grained routing within the selected pool. This hierarchical routing enables efficient sharing of parameters across related medical tasks (e.g., diagnosis coding, outcome prediction, report generation) while maintaining task-specific specialization. The dual-gate design reduces the total number of trainable parameters by 85\% compared to full MoE fine-tuning while maintaining comparable performance.

The MoELoRA approach \cite{rai2024moe_lora} further advances this direction by incorporating contrastive learning to guide expert specialization during training. By encouraging experts to develop maximally distinct representations, MoELoRA mitigates the expert collapse problem that is particularly acute when the number of experts exceeds the natural number of data clusters in medical corpora.

\subsection*{Clinical text understanding and information extraction}

Beyond LLMs, MoE architectures have been applied to specialized clinical NLP tasks. Clinical text contains complex medical terminology, abbreviations, and implicit clinical reasoning that varies across specialties. MoE-based models have shown improved performance in clinical named entity recognition by routing text from different medical domains through domain-specialized experts. Similarly, in clinical relation extraction, different expert pathways can specialize in different relation types (drug-disease, symptom-diagnosis, procedure-outcome), improving extraction accuracy for rare relation types that are underrepresented in training data.

\subsection*{Reducing hallucination in medical LLMs}

A critical challenge for medical LLMs is the generation of factually incorrect or clinically dangerous outputs (hallucinations). MoE architectures offer a unique advantage in hallucination mitigation through expert-level confidence estimation and routing-based quality control. When the gating network exhibits low confidence in expert selection for a given input, this uncertainty signal can be used to trigger abstention or human review, providing a built-in safety mechanism that is absent in dense models. \cite{bechar2025biocancer_moe} demonstrated that knowledge-guided routing can reduce hallucination rates by encoding domain ontologies as constraints on the gating function, preventing the model from generating outputs that contradict established medical knowledge.

The simulation of multi-disciplinary team reasoning through MoE \cite{chen2026cor_mole} provides another avenue for hallucination reduction: by requiring consensus among multiple expert perspectives, the model is less likely to produce outlier outputs that arise from overconfident predictions of a single expert pathway. This collaborative verification mechanism parallels clinical practice, where diagnostic errors are reduced through peer consultation and second opinions.
